\section{Đạo hàm đi qua ô nhớ}
\label{sec:ch04-dao-ham}

\index{vòng quay lỗi không đổi!đạo hàm qua nó}%
Chương~\ref{ch:phan-tich} đo \(\norm{\jac{\vect{a}_t}{\vect{a}_k}}\) và thấy
nó tắt theo cấp số nhân. Câu hỏi tương ứng ở đây là
\(\jac{\vect{c}_t}{\vect{c}_{t-1}}\) bằng bao nhiêu. Lấy đạo hàm
của~\eqref{eq:ch04-cap-nhat-c}:

\begin{equation}
  \jac{\vect{c}_t}{\vect{c}_{t-1}}
    = \mat{I}
    + \jac{\left(\vect{i}_t \odot g_{\mathrm{in}}(\vect{a}^{c}_t)\right)}{\vect{c}_{t-1}}.
  \label{eq:ch04-jac-tach}
\end{equation}

Số hạng thứ nhất là \(\mat{I}\) vì \(\vect{c}_{t-1}\) đứng trong phương trình
với hệ số 1. Số hạng thứ hai là chỗ cần cẩn thận, vì thoạt nhìn nó bằng không:
\(\vect{i}_t\) và \(\vect{a}^{c}_t\) là hàm của \(\vect{x}_t\) và
\(\vect{h}_{t-1}\), không phải của \(\vect{c}_{t-1}\). Nhưng
\(\vect{h}_{t-1} = \vect{o}_{t-1} \odot g_{\mathrm{out}}(\vect{c}_{t-1})\), nên
nó \emph{là} hàm của \(\vect{c}_{t-1}\), và số hạng thứ hai không bằng không.

Có một đường vòng: tín hiệu rời khỏi cell qua \(g_{\mathrm{out}}\), bị
\tn{cổng ra}{output gate} co lại, rồi quay vào cell ở bước sau qua
\tn{cổng vào}{input gate} và qua cell input. Viết ra đầy đủ, với
\(\mat{D} = \diag(\vect{o}_{t-1} \odot g_{\mathrm{out}}'(\vect{c}_{t-1}))\):

\begin{equation}
  \jac{\vect{c}_t}{\vect{c}_{t-1}}
    = \mat{I}
    + \Big[
        \diag\!\left(g_{\mathrm{in}}(\vect{a}^{c}_t)\right)
        \diag\!\left(\sigma'(\vect{a}^{i}_t)\right) \mat{W}_{hi}
      + \diag(\vect{i}_t)
        \diag\!\left(g_{\mathrm{in}}'(\vect{a}^{c}_t)\right) \mat{W}_{hc}
      \Big] \mat{D}.
  \label{eq:ch04-jac-day-du}
\end{equation}

Một chỗ dễ sai khi tự dẫn lại: \(\vect{o}_{t-1}\) trong \(\mat{D}\) là cổng ra
của bước \emph{trước}, và nó là hàm của \(\vect{h}_{t-2}\) chứ không phải của
\(\vect{c}_{t-1}\). Coi nó là hàm của \(\vect{c}_{t-1}\) thì
thêm một số hạng không có thật. Tôi kiểm~\eqref{eq:ch04-jac-day-du} bằng máy,
so với \code{torch.autograd.functional.jacobian}, và lệch lớn nhất dưới
\(10^{-12}\); test khóa lại chuyện đó ở tag \code{ch04}.

\subsection*{Phép cắt làm số hạng thứ hai biến mất}

\index{lan truyền ngược cắt cụt!hệ quả cho CEC}%
Bài không tính~\eqref{eq:ch04-jac-day-du}. Thuật toán của bài dùng
\tn{lan truyền ngược cắt cụt}{truncated backpropagation}, và phép cắt ấy thay
đúng ba đạo hàm bằng không. Hai trong ba đạo hàm ấy là hai đường mà số hạng thứ
hai đi qua, tức \(\mat{W}_{hi}\) và \(\mat{W}_{hc}\)
trong~\eqref{eq:ch04-jac-day-du}; đường thứ ba đi qua cổng ra và không có mặt ở
đây, vì \(\vect{o}_t\) không xuất hiện trong \(\vect{c}_t\).
Mục~\ref{sec:ch04-cat-gradient} liệt kê cả ba; ở đây chỉ cần hệ quả:

\begin{equation}
  \jac{\vect{c}_t}{\vect{c}_{t-1}} \overset{\text{cắt}}{=} \mat{I}.
  \label{eq:ch04-jac-cat}
\end{equation}

Ma trận đơn vị. Không phải một ma trận gần ma trận đơn vị. Và vì tích của các
ma trận đơn vị vẫn là ma trận đơn vị,

\begin{equation}
  \jac{\vect{c}_t}{\vect{c}_k} = \mat{I} \quad \text{với mọi } k \le t.
  \label{eq:ch04-jac-moi-khoang-cach}
\end{equation}

Đây là câu trả lời cho câu hỏi chương~\ref{ch:phan-tich} để lại. Thừa số không
những gần 1, nó bằng 1, ở mọi bước, ở mọi khoảng cách, không cần huấn luyện giữ
nó ở đó.

\subsection*{Ba đường cong, một trục}

Tôi đo cả ba trên cùng một trục khoảng cách. Dòng
\tn{mạng hồi quy}{recurrent neural network} thường, viết tắt là RNN, dùng đúng
ma trận, đúng seed và đúng \tn{bán kính phổ}{spectral radius} 0.9 mà
chương~\ref{ch:phan-tich} dùng, nên nó phải trùng khít số chương đó in ra, và
có một test khóa lại chuyện đó.

\begin{measured}{đạo hàm qua CEC so với đạo hàm qua một RNN thường}
64 đơn vị ẩn, 100 bước. Trọng số \tn{ô nhớ}{memory cell} lấy trong
\([-0.1, 0.1]\), đúng khoảng bài dùng cho các thí nghiệm 3 tới 6:

\begin{minted}{text}
model            l=1         l=10        l=50        l=100
plain RNN 0.9    9.0000e-01  3.4200e-01  5.0259e-03  2.5902e-05
CEC truncated    1.0000e+00  1.0000e+00  1.0000e+00  1.0000e+00
CEC full w=0.1   1.0000e+00  2.0016e+00  2.9078e+01  1.2834e+02
\end{minted}

Dòng giữa không phải làm tròn. Sai lệch lớn nhất so với 1.0 trên cả 100 bước,
trong float64, là 0.0e+00.

Dòng dưới là dòng tôi không đoán trước. Bỏ phép cắt đi thì đạo hàm không tắt,
nó \emph{tăng}, và tăng ngay ở thang khởi tạo của chính bài báo: 29.078 ở
khoảng cách 50 và 128.34 ở khoảng cách 100. Chạy lại với trọng số lớn hơn mười
lần thì con số ở khoảng cách 100 lên 3629.8, nên đây không phải hiệu ứng của
một lần rút may rủi.

Đo bằng \code{experiments/ch04\_flow.py} tại tag \code{ch04}.
\end{measured}

Hình~\ref{fig:ch04-dong-loi} vẽ ba đường trên thang log.

\begin{figure}[htbp]
  \centering
  % Chuẩn của tích đạo hàm theo khoảng cách, ba đường trên cùng một trục.
% Số liệu: research/2026-08-ch04-lstm.md, mục "Đo 3". Chỉ bốn khoảng cách được
% đo (l = 1, 10, 50, 100) cộng với l = 0; các đoạn giữa là nội suy thẳng.
% Toạ độ: x_cm = l * 0.1; y_cm = (log10(norm) + 5) * 0.8.
% Mono-safe: nét liền, nét đứt và nét chấm phân biệt ba đường.
\begin{tikzpicture}[
  >= Stealth,
  truc/.style   = {->, >={Stealth[length=2mm]}, gray!70, line width = 0.5pt},
  rnn/.style    = {line width = 0.9pt, black, densely dotted},
  cec/.style    = {line width = 1.1pt, black},
  day/.style    = {line width = 0.9pt, black!70, densely dashed},
  ghichu/.style = {font = \footnotesize},
  nho/.style    = {font = \scriptsize},
  diem/.style   = {circle, fill = black, inner sep = 0pt, minimum size = 1.6mm},
]

  % --- trục
  \draw[truc] (-0.3, 0) -- (11.0, 0)
    node[right, ghichu, black] {\(l = t - k\)};
  \draw[truc] (0, -0.3) -- (0, 6.3)
    node[above, ghichu, black] {\(\norm{\partial \cdot_t / \partial \cdot_k}\)};

  % --- vạch trục x
  \foreach \lx/\llab in {0/0, 1/10, 5/50, 10/100} {
    \draw[gray!70] (\lx, 0) -- (\lx, -0.12);
    \node[below, nho] at (\lx, -0.12) {\llab};
  }
  % --- vạch trục y, thang log
  \foreach \ly/\llab in {0.8/{\(10^{-4}\)}, 2.4/{\(10^{-2}\)},
                         4.0/{\(1\)}, 5.6/{\(10^{2}\)}} {
    \draw[gray!70] (0, \ly) -- (-0.12, \ly);
    \node[left, nho] at (-0.12, \ly) {\llab};
  }
  % --- đường ngang tại 1, để mắt bắt được mức "không đổi"
  \draw[gray!35, line width = 0.4pt] (0, 4.0) -- (10.6, 4.0);

  % --- CEC có cắt gradient: bằng 1 ở mọi khoảng cách
  \draw[cec] (0, 4.0) -- (10.0, 4.0);
  \node[ghichu, anchor = south west] at (6.2, 4.06) {CEC, có cắt gradient};

  % --- CEC không cắt: 1.0, 2.0016, 29.078, 128.34
  \draw[day] (0, 4.0) -- (1.0, 4.241) -- (5.0, 5.171) -- (10.0, 5.687);
  \foreach \px/\py in {1.0/4.241, 5.0/5.171, 10.0/5.687} {
    \node[diem] at (\px, \py) {};
  }
  \node[ghichu, anchor = south east, align = right] at (9.9, 5.75)
    {CEC, không cắt gradient};

  % --- plain RNN, bán kính phổ 0.9: 0.9, 0.342, 5.0259e-3, 2.5902e-5
  \draw[rnn] (0, 4.0) -- (0.1, 3.963) -- (1.0, 3.627) -- (5.0, 2.161)
             -- (10.0, 0.331);
  \foreach \px/\py in {1.0/3.627, 5.0/2.161, 10.0/0.331} {
    \node[diem] at (\px, \py) {};
  }
  \node[ghichu, anchor = north east, align = right] at (9.9, 1.15)
    {RNN thường, \(\rho(\mat{W}_{hh}) = 0.9\)\\ (số của chương~3)};

  % --- chấm là chỗ đo thật
  \node[nho, anchor = north west, black!60, align = left] at (0.15, 6.25)
    {chấm tròn = khoảng cách thật sự được đo;\\
     đoạn nối giữa hai chấm là nội suy};

\end{tikzpicture}

  \caption{Chuẩn của tích đạo hàm theo khoảng cách \(l = t - k\), thang log.
    Đường liền nằm ngang là CEC có cắt gradient, bằng 1 ở mọi khoảng cách.
    Đường chấm đi xuống là RNN thường của chương~\ref{ch:phan-tich} với
    \(\rho(\mat{W}_{hh}) = 0.9\). Đường đứt đi lên là cùng ô nhớ ấy khi bỏ phép
    cắt: đường rò rỉ ra khỏi cell rồi quay lại không tắt mà phình theo khoảng
    cách. Bốn khoảng cách được đo là \(l = 1, 10, 50\) và \(100\); chấm tròn
    đánh dấu ba cái sau, vì ở \(l = 1\) cả ba đường còn nằm chồng lên nhau.}
  \label{fig:ch04-dong-loi}
\end{figure}

\subsection*{Tín hiệu bị co ở hai đầu, không co ở giữa}

\index{vòng quay lỗi không đổi!chỉ co ở hai đầu}%
Phương trình~\eqref{eq:ch04-jac-moi-khoang-cach} nói về đường bên trong CEC.
Tín hiệu lỗi thì phải đi vào rồi đi ra, và hai chỗ ấy không miễn phí. Phụ lục
A.2 của bài viết công thức đầy đủ, và câu kết luận ngay dưới nó đáng nhớ
nguyên văn: dòng lỗi chỉ bị co tại thời điểm \(t\), lúc nó vào cell, và tại
\(t - q\), lúc nó ra, \enquote{but not in between}.

Cụ thể, vào thì nhân với \(\vect{o}_t \odot g_{\mathrm{out}}'(\vect{c}_t)\),
tức qua cổng ra và qua đạo hàm của hàm nén output; ra thì nhân với một thừa số
có \(\vect{i}_{t-q}\) và \(g_{\mathrm{in}}'\) trong đó. Cả hai thừa số nhỏ hơn
1, nên tín hiệu \emph{có} yếu đi. Điều nó không làm là yếu đi theo \(q\). Hai
lần co, bất kể quãng đường dài bao nhiêu, thay cho \(q\) lần co.

Đó là toàn bộ khác biệt giữa một hằng số và một hàm mũ, và nó là lý do bài dám
nói tới quãng vượt một nghìn bước.
